\section{Policy Implications}
\label{sec:policy-implications}

\subsection{IIJA Resilience Funding Is Insufficient for the Priority Portfolio}

The Infrastructure Investment and Jobs Act \citep{IIJA2021} allocated \$3.75 billion over five years for highway resilience projects under the PROTECT (Promoting Resilient Operations for Transformative, Efficient, and Cost-Saving Transportation) formula and discretionary programs. This represents a significant federal resilience investment --- but it is insufficient for Phase 1 alone. The Donner Pass freight tunnel (\$4.0B) exceeds the total five-year PROTECT program authorization by \$250 million, and requires a dedicated authorization outside the formula program structure.

This is not a reason to deprioritize the Donner tunnel; it is a reason to pursue dedicated authorization. The precedent exists: the Portal Bridge replacement on Amtrak's Northeast Corridor received a \$766M FRA grant with a total project cost of \$3.3B, funded through a combination of FRA discretionary grants, state contributions, and Gateway Program authority. A similar dedicated authorization structure for the Donner freight tunnel --- Federal contribution 60\%, FHWA freight formula contribution 20\%, California and Nevada state contributions 20\% --- is administratively feasible under existing statutory authority.

\subsection{FHWA Resilience Grant Allocation Methodology}

Current FHWA climate resilience grant prioritization under PROTECT uses criteria that are primarily D1-equivalent: FEMA Special Flood Hazard Area miles, historical severe weather event frequency, and projected climate hazard intensity. These criteria correctly identify corridors with high environmental exposure. They do not account for network topology --- specifically, whether the at-risk corridor is the only path for a significant share of national freight or one of several adequate paths.

This omission has a systematic consequence: PROTECT grants are currently allocated toward corridors with high D1 scores regardless of their B1 isolation level. A corridor with D1 = 8.0 and B1 = 2.0 (I-95 South Carolina --- high flood exposure, good alternates) competes on equal D1 terms with a corridor with D1 = 7.8 and B1 = 8.3 (Donner Pass), despite the fact that Donner closures produce 20 times the national freight disruption per event.

The recommendation is to incorporate the B1$\times$D1 compound priority score as a required criterion in PROTECT discretionary grant evaluation. Specifically: any corridor applying for PROTECT resilience investment funding should be required to submit its ROUTE corpus B1 and D1 scores (or equivalent alternate-route-adequacy and climate-exposure metrics), and the compound score should be weighted at no less than 30\% of total application score. This change does not require statutory amendment; FHWA has authority to define grant evaluation criteria under existing IIJA Title 23 authority.

\subsection{National Critical Highway Infrastructure Designation}

The top four compound exposure corridors --- Donner Pass, Gulf Coast I-10 Louisiana, Snoqualmie Pass, and Siskiyou Pass --- should be designated as National Critical Highway Infrastructure under the FAST Act's Critical Transportation Infrastructure Protection authority (49 U.S.C. \S 70101 et seq.). This designation triggers four federal authorities not available to standard NHS corridors:

\begin{enumerate}
  \item Mandatory federal funding participation above normal NHS apportionment formula (up to 100\% federal share for critical resilience projects)
  \item Access to DoD Infrastructure Protection funds for corridors with B4 scores above 6.0 (Donner and Snoqualmie serve Pacific military logistics)
  \item Expedited environmental review under the One Federal Decision framework for critical infrastructure
  \item Mandatory state reporting on closure events, durations, and detour costs to FHWA --- creating the data infrastructure for ongoing compound-exposure monitoring
\end{enumerate}

\subsection{The 2050 Data Gap and Priority Order Instability}

D1 scores in this analysis use 2024 baseline climate data. NOAA 2050 intermediate sea-level-rise projections and IPCC RCP 4.5 precipitation-intensity projections are available but have not been integrated into the ROUTE corpus D1 scores. This data gap has a material consequence for investment sequencing.

Under 2050 climate projections, the Gulf Coast I-10 Louisiana D1 score is projected to rise from 8.4 to approximately 9.1, driven by a 0.5--1.0 meter sea-level-rise increase in storm surge reach. This would shift the Gulf Coast compound score from 57.1 to approximately 61.9 (6.8 $\times$ 9.1), narrowing the gap with Donner (64.7) substantially. Under high-range SLR scenarios (1.0m by 2050), Gulf Coast I-10 Louisiana could move above Donner in the compound priority ranking by 2035 --- reversing the Phase 1/Phase 2 sequencing.

The policy implication is that the Phase 2 Gulf Coast investment should not wait for 2050 SLR confirmation. The current 2024 data already supports second-priority investment; the trajectory data suggests the gap between Gulf Coast and Donner will narrow over the investment horizon. Beginning Gulf Coast environmental review and design work concurrently with Phase 1 Donner tunnel authorization is consistent with both the current priority ordering and its likely 2050 evolution.

\subsection{Complementarity with Interstate 2.0 Managed Freight Lanes}

The compound exposure investment program is complementary to, not substitutive of, the Interstate 2.0 managed freight lane program \citep{ROUTE_D1}. Managed freight lanes address throughput and reliability under normal operating conditions: they increase capacity, enforce truck-only separation, and enable time-guaranteed freight delivery slots. They do not address compound exposure, because managed freight lanes still close during the 72-hour Donner blizzard, the Gulf Coast hurricane surge, and the Snoqualmie avalanche cycle. A managed freight lane through Donner Pass is closed for the same 50 annual events as the general traffic lanes; the tunnel replaces the exposure entirely.

The investment implication is additive: compound-exposure hardening (B3 resilience holes portfolio) and managed freight lanes (I2.0 program) address different failure modes and serve different operational objectives. A corridor that receives both investments --- for example, I-80 with a Donner freight tunnel plus managed freight lanes in the Sacramento Valley and Salt Lake City segments --- is both resilient to climate closure and operationally efficient under normal conditions. Neither investment makes the other redundant.
